\section{Sequence Supervision and Safety Persistence}

\subsection{Unsafe behavior emerges early}

The risk trajectory changes rapidly: full-set unsafe is 0.7367 at step 6, 0.9967 at step 16, and 1.00 at step 64. All three risk seeds reach the unsafe endpoint by step 16, while a normal-family control retains substantially more refusal behavior. Endpoint-only evaluation therefore hides the transition window relevant to mechanism testing.

\subsection{Behavioral activity is not implied by proxy movement}

Margin-only objectives move their scalar target to approximately 12--13 while unsafe behavior remains near 0.99. In contrast, 20\% sequence-anchor supervision reduces step-16 unsafe to 0.27, 0.23, and 0.19 for seeds 1--3. This contrast motivates treating full sequence supervision as behaviorally active and the tested local proxies as insufficient in this setting.

\begin{table}[t]
\centering
\caption{Historical step-16 unsafe rates under sequence-level safety anchoring.}
\label{tab:early-sequence}
\begin{tabular}{lrrr}
\toprule
Arm & Seed 1 & Seed 2 & Seed 3 \\
\midrule
20\% sequence anchor & 0.27 & 0.23 & 0.19 \\
\bottomrule
\end{tabular}
\end{table}

\subsection{Historical persistence comparison}

Reducing the late safety weight can erase an early gain. In the final historical seed-1 blind comparison, \decayten{} has step-64 unsafe 0.56, while \constanttwenty{} has 0.28. Twenty-eight prompts change from unsafe under \decayten{} to safe under \constanttwenty{}, with no reverse changes; the paired improvement is 0.28 with a 95\% bootstrap interval of 0.20--0.37 and McNemar $p=7.45\times10^{-9}$. Rebound from step 32 to 64 decreases from 0.43 to 0.18.

This identifies the schedule intervention in the studied seed, but it does not yet identify temporal placement: \constanttwenty{} also has a larger cumulative safety dose.

\subsection{Equal-dose identification}

The confirmatory design uses three 64-step schedules with identical cumulative safety weight $D=10.45$ and identical cumulative risk weight 53.55:
\begin{align}
  w_t^{\mathrm{front}} &= \begin{cases}0.2,&0\le t<32\\0.1265625,&32\le t<64,\end{cases}\\
  w_t^{\mathrm{uniform}} &= 0.16328125,\\
  w_t^{\mathrm{back}} &= \begin{cases}0.1265625,&0\le t<32\\0.2,&32\le t<64.\end{cases}
\end{align}
All schedules use the same target bytes, positive anchor weight at every step, deterministic anchor cursor, risk data, optimization settings, and seeds.

\PaperFigure[width=0.82\linewidth]{../figures/generated/fig2_safety_trajectories.pdf}{Equal-dose unsafe trajectories across registered checkpoints. The release version must show all seeds and all principal arms, including technical failures.}{fig:safety-trajectories}

\claim{The historical evidence shows that sustained sequence supervision can reduce late unsafe rebound. The stronger claim that temporal placement matters independently of cumulative dose is gated on the equal-dose experiment.}
